\section{Related Work}
\label{sec:related-work}

\paragraph{Evolutionary NAS with surrogates.} NSGA-Net searches architectures with NSGA-II, training every
candidate fully~\citep{lu2019nsganet}. NSGANetV2 adds an accuracy predictor refined online and inherits weights
from a supernet, cutting the number of fully evaluated architectures to 350~\citep{lu2020nsganetv2}, at the
documented risk of rank disagreement between weight-shared and stand-alone accuracy~\citep{yu2020evaluatingnas}.
White et al.\ identify multi-objective evolutionary search as the most common approach to multi-objective NAS,
with NSGANetV2 as an example~\citep{white2023insights}; Liu et al.\ list learned surrogates as one of several
cost-reduction techniques alongside weight sharing and zero-cost proxies~\citep{liu2023survey}. Later work in
this lineage moves toward ranking-based predictors~\citep{lu2022mosegnas}. None of three recent surveys lists
linkage learning as a distinct NAS method category~\citep{liu2023survey,white2023insights,ozcelik2026survey}.

\paragraph{Linkage learning and P3.} GOMEA and P3 build a linkage tree from statistical dependencies in the
population and use it to guide crossover, so that groups of dependent variables are not
disrupted~\citep{thierens2011optimal,goldman2014parameterless}. P3 replaces a single fixed-size population with
a pyramid of populations of growing size (Figure~\ref{fig:p3-pyramid}), adding a level only once existing levels
stop improving, which is why no population size is chosen in advance. This structure is central to the
diagnosis in Section~\ref{sec:results}. Unlike the gray-box setting that the GOMEA line often
exploits~\citep{whitley2016graybox,bouter2023library}, NAS accuracy is not decomposable into cheaply recomputable
subfunctions, so here every variation step must be verified by a full evaluation or screened by a surrogate.

\begin{figure}[t]
    \centering
    \begin{tikzpicture}[
        lvl/.style={draw, rounded corners, minimum height=6mm, font=\scriptsize, align=center},
    ]
        \node[lvl, minimum width=14mm] (l1) at (0,0)     {level 1};
        \node[lvl, minimum width=22mm] (l2) at (0,-9mm)  {level 2};
        \node[lvl, minimum width=30mm] (l3) at (0,-18mm) {level 3};
        \node[lvl, minimum width=38mm] (l4) at (0,-27mm) {level 4};
        \draw[{Stealth[length=2mm]}-, dashed] (l4.south) -- ++(0,-6mm)
            node[below, font=\tiny, align=center, text width=32mm]
            {new level added only once existing levels stop yielding improved solutions};
    \end{tikzpicture}
    \caption{P3's population pyramid: each level holds a population of growing size; a new level is added only
    when the levels below it stop producing improved solutions.}
    \Description{A vertical stack of four boxes labelled level 1 through level 4, growing wider from top to
    bottom, with an arrow noting that a new level is added only once existing levels stop yielding improved
    solutions.}
    \label{fig:p3-pyramid}
\end{figure}

\paragraph{Surrogate-assisted linkage learning.} CS-GOMEA pairs GOMEA with a convolutional surrogate on
synthetic combinatorial problems, and flags as open the surrogate's failure to use linkage information in its own
construction~\citep{dushatskiy2019csgomea}. Dushatskiy et al.\ integrate a surrogate with a P3 variant on an
ensemble-learning problem, where an SVR outperformed a neural
predictor~\citep{dushatskiy2021novelsurrogateassistedevolutionaryalgorithm}. LyMPuS integrates a surrogate
directly with linkage discovery and validates it with P3 on NK landscapes, Ising spin glasses, and
Max3Sat~\citep{przewozniczek2026lympus}; its empirical variant, eLyMPuS, discovers the dependency structure
incrementally rather than assuming it known. P3Net's surrogate follows eLyMPuS in being relative and
linkage-aware over an incompletely discovered structure, but regresses a continuous fitness difference rather
than making a discrete better/worse/ambiguous comparison, which places its mechanism closer to CS-GOMEA's and
forgoes eLyMPuS's monotonicity-conditional recovery guarantee.

\paragraph{Scale dependence.} The lineage's advantage over structure-blind search is documented to grow with
problem size. Goldman and Punch report a lower order of evaluation complexity for P3 than five competitors over
two orders of magnitude of problem size~\citep{goldman2015fast}; Qiao and Gallagher derive a GOMEA runtime bound
on concatenated traps whose gap over a $(1{+}1)$~EA widens with both subfunction count and
length~\citep{qiao2024gomearuntime}. Within NAS, den Ottelander et al.\ find MO-GOMEA, NSGA-II, local search, and
random search indistinguishable on NAS-Bench-101, while MO-GOMEA does beat random search on their larger MacroNAS
benchmarks~\citep{denottelander2021localsearch}.

\paragraph{Closest prior work.} Bartnik adapts SA-P3-GOMEA to multi-objective NAS on NAS-Bench-201, optimising
validation accuracy against measured GPU energy~\citep{bartnik2026evolutionary}. The surrogate is an absolute
regressor over the flat genotype (SVR, MLP, random forest, or gradient boosting); baselines are MO-LS and
surrogate-free MO-P3-GOMEA; the genotype is architecture-only. P3Net differs on all three axes: a relative,
linkage-aware surrogate; NSGA-II and NSGANetV2 as baselines; and a joint architecture-and-hyperparameter genotype,
to which, to our knowledge, no linkage-learning EA has previously been applied. Tran et al.\ also apply GOMEA to
NAS, but reduce cost through a training-free metric rather than a learned surrogate~\citep{tran2023gomeanas}.

\paragraph{LLM-based search.} Large language models enter the same search loop in three roles. As the optimiser
itself, GPT-4 proposes NAS-Bench-201 cells over ten iterations~\citep{zheng2023genius}, and a code-generating LLM
serves as the variation operator of a quality-diversity NAS~\citep{nasir2024llmatic}; outside NAS, LLM-driven
evolutionary search is competitive on the TSP only up to 20 nodes~\citep{liu2024lmea}, and LLMs underperform on
purely numerical black-box problems~\citep{huang2024llmblackbox}. As a surrogate, in-context LLM predictors match
conventional surrogates on 5- and 10-dimensional test functions at a higher inference
cost~\citep{hao2024llmsurrogate}, and an LLM performance predictor for machine-translation NAS attains the lowest absolute error but a
slightly weaker rank correlation than its baseline predictor~\citep{jawahar2024llmpp} -- ranking being what a
relative surrogate is used for. As a prior, LLMs warm-start or stand in for Bayesian
optimisation~\citep{liu2024llambo,zhang2024llmhpo} and prune a NAS search space by transferring design principles
from previously searched tasks~\citep{zhou2025lapt}. The documented gains concentrate in the early, sparse-data
phase~\citep{liu2024llambo}, on spaces of two to five hyperparameters and 10--30 evaluations in the main results of
Zhang et al.~\citep{zhang2024llmhpo}. Several studies trace them to prior knowledge rather than to feedback: LLMs
help Bayesian optimisation over molecules only when pretrained or fine-tuned on domain
data~\citep{kristiadi2024sober}, replacing observed outcomes with permuted labels leaves LLM-based experimental
design unaffected~\citep{gupta2025llmbo}, and within a fixed search space CMA-ES and TPE outperform LLM agents, a
hybrid of the two doing best~\citep{ferreira2026centaur}. Linkage learning is the opposite case: it uses no prior
and learns structure only from evaluations, with an advantage that grows with problem size (see
\emph{Scale dependence}). We found no direct comparison of the two families. We include no LLM-based arm, for
three reasons: it would draw on information outside the evaluation budget to which every other arm is held;
NAS-Bench-201 is public, and the GENIUS authors cannot rule out that GPT-4 has seen the benchmarks they
use~\citep{zheng2023genius}; and our question is why a linkage-learning engine fails to separate from random
search, which an added arm would not answer. The best-documented LLM role, warm-starting the sparse early phase,
nonetheless targets the uniform-random bootstrap that our diagnosis implicates (Section~\ref{sec:results}).

\paragraph{Joint NAS+HPO benchmarks and their limits.} JAHS-Bench-201 answers queries through a learned surrogate
over a mixed categorical/continuous joint space~\citep{bansal2022jahsbench}; NAS-HPO-Bench-II is a fixed lookup
table over architecture, learning rate, and batch size~\citep{hirose2021nashpobenchii}. Only the latter admits an
exactly enumerable oracle Pareto front, which we exploit for IGD+ (Section~\ref{sec:results}). Guerrero-Viu et
al.\ establish SH-EMOA and MO-BOHB as baselines for this joint setting~\citep{guerreroviu2021bagofbaselines}, none
of them a linkage-learning EA. Both benchmarks share the NAS-Bench-201 cell space~\citep{dong2020nasbench201} and
small-scale image classification. The benchmarking literature documents that search-space design accounts for
much of a method's reported performance and that rankings transfer poorly across
benchmarks~\citep{yang2020nasfrustratinglyhard,mehta2022nasbenchsuite}, that this particular cell space is
redundant and skewed toward already-good architectures~\citep{lopes2023nasbenchdesign}, and that conclusions
transfer inconsistently across task domains~\citep{tu2022nasbench360,duan2021transnasbench101}. Any result on
this family should be read against that.
